% Anexo GENERADO por herramientas/docx_a_anexo.py desde 2026-11-19_yo_EETT-estacion-experimental.docx
% Para corregir el contenido, editar el Word original y volver a generar.

\section{Sector 1: recolección e impulsión}

\subsection{Recolección e impulsión de aguas servidas}

\subsubsection*{Suministro e Instalación de Fosas Sépticas (3.250 L y 1.250 L)}

Esta partida comprende el suministro, transporte, excavación, colocación, nivelación, relleno y prueba de hermeticidad de dos fosas sépticas estancas de polietileno de alta densidad (PEAD). Se utilizará un estanque de 3.250 Litros para la recepción general y uno de 1.250 Litros para el ajuste secundario, destinados al tratamiento primario mediante digestión anaerobia previa al envío hacia los humedales. La ejecución cumplirá con la NCh2038, el D.S. N° 236/26 del MINSAL y los manuales del fabricante. Los estanques serán monobloque con protección UV, equipados con tapas de registro roscadas y conexiones sanitarias en PVC Ø 110 mm con tees deflectoras de natas. La instalación requiere una excavación con holgura perimetral de 30 a 50 cm y una base de arena limpia de 15 cm. Durante el relleno lateral compactado por capas con arena o estabilizado fino, los estanques se llenarán simultáneamente con agua limpia para equilibrar presiones y evitar deformaciones estructurales, asegurando una pendiente de llegada no inferior al 2\%. Antes de cubrir superiormente el estanque, se realizará una prueba de hermeticidad hidrostática manteniendo la estructura llena durante 24 horas continuas para verificar la ausencia total de filtraciones. La unidad de medida para esta partida es por \textbf{Unidad (un)}.

\subsubsection*{Cámaras de Inspección Sanitarias}

Consiste en la construcción o suministro de cámaras de inspección domiciliarias destinadas al registro, cambios de dirección y mantenimiento de la red de alcantarillado de las instalaciones residenciales y de oficina. Se ajustarán a la NCh2592 y al RIDAA. Podrán utilizarse cámaras prefabricadas de PEAD Ø 600 mm o, equipadas con marco, tapa estanca y banquetas con pendiente del 20\% hacia la canaleta central. Se asentarán sobre una base de hormigón pobre o arena estabilizada de 10 cm y se acoplarán a las tuberías de PVC mediante juntas elásticas o soldadura química. El control de calidad exigirá la verificación visual de pendientes, acabado de banquetas e inspección de hermeticidad en las uniones. La unidad de medida es por \textbf{Unidad (un)}.

\subsubsection*{Suministro e Instalación de Tubería Sanitaria PVC 110 mm}

Abarca el suministro e instalación de la canalización de PVC Sanitario Ø 110 mm para la conducción por gravedad del agua servida bruta. Se regirá por las normas NCh1635 y NCh2282/2. Las tuberías serán de PVC rígido sanitario de serie normal, unidas mediante junta elástica o adhesivo químico NCh2252 previa limpieza de extremos. El tendido se ejecutará en zanja asegurando una pendiente continua entre el 1\% y el 3\%, apoyando el tubo sobre una cama de arena limpia de 10 cm de espesor. La unidad de medida es por \textbf{Metro lineal (ml)}.

\subsubsection*{Suministro e Instalación de PVC Hidráulico 50 mm (Impulsión)}

Comprende el suministro y tendido de tubería de PVC Hidráulico de 50 mm de diámetro exterior, Clase 10 (PN10), para la línea de impulsión a presión del agua servida pretratada. Se ejecutará bajo los estándares de la norma NCh399. La tubería se instalará en zanja sobre cama de arena de 10 cm e incorporará una cinta de advertencia subterránea antes del tapado. A zanja descubierta. La unidad de medida es por \textbf{Metro lineal (ml)}.

\subsubsection*{Suministro e Instalación de Fitting}

Incluye la totalidad de las piezas especiales de unión, codos de 45° y 90°, tees, reducciones, coplas y/o curvas necesarias para interconectar los nudos del Sector 1.1. Todas las piezas serán inyectadas en PVC sanitario o hidráulico PN10 de primera calidad y deberán montarse asegurando hermeticidad mediante soldadura química o juntas elásticas mecánicas. Su aprobación quedará supeditada a las pruebas de presión e impermeabilidad de las redes principales. Se medirá de forma \textbf{Global (gl)}.

\subsubsection*{Excavaciones y Cama de Arena (Sector 1.1)}

Comprende los movimientos de tierra requeridos para zanjas y fosas, así como la colocación de la cuna de apoyo en arena limpia. Las excavaciones se realizarán cumpliendo con la NCh349 sobre seguridad en excavaciones. Se empleará arena limpia de río o mina, exenta de materia orgánica y piedras angulosas, dispuesta en capas de 10 a 15 cm bajo las tuberías y como relleno de protección lateral. Se inspeccionarán las cotas de sello y la compactación de la cuna previa al tendido. La unidad de medida es por \textbf{Metro cúbico (m³)}.

\subsection{Recolección e impulsión de aguas del estero}

\subsubsection*{Sistema de Captación (Válvula de Pie y Jaula de Contención)}

Abarca el suministro e instalación de la obra de toma en el estero, constituida por una válvula de pie Ø 50 mm y una jaula metálica de protección estructural. La válvula contará con un canastillo filtrante de acero inoxidable AISI 316, mientras que la jaula se fabricará con perfiles de acero galvanizado. La jaula se anclará firmemente al lecho o ribera del estero, manteniendo la válvula suspendida a no menos de 20 cm del fondo para evitar la succión de sedimentos. Se evaluará la retención de columna de agua y la ausencia de cavitación durante la operación de la bomba. Se medirá por \textbf{Unidad (un)}.

\subsubsection*{Grava de Protección (64 mm)}

Consiste en la colocación de un manto de grava gruesa de 64 mm alrededor de la jaula de toma para actuar como disipador y filtro grueso mecánico. El material responderá a la norma NCh163, utilizándose grava lavada limpia. Se dispondrá manualmente asegurando un recubrimiento homogéneo perimetral. La unidad de medida es por \textbf{Metro cúbico (m³)}.

\subsubsection*{Caseta y Tablero Eléctrico Sector Estero}

Comprende la ejecución de una losa de fundación de hormigón G10 de 10 cm de espesor y el suministro/montaje de un gabinete estanco IP65 para alojar el equipo de bombeo y su tablero de control. El tablero eléctrico se configurará para energización posterior, con instalación a 1,5 m. El gabinete se fijará sólidamente a la losa y se cableará a través de canaletas ranuradas con peinetas identificadoras. Se realizarán ensayos de aislamiento, continuidad y pruebas funcionales de protecciones. Se medirá por \textbf{Unidad (un)}.

\subsubsection*{Red de Conducción, Válvulas, Fitting y Tierras Sector Estero}

Abarca la instalación de la red en PVC Hidráulico Ø 25 mm PN10, cuatro válvulas de bola/mariposa de 25 mm con unión americana, sus fitting inyectados y las obras de tierra asociadas (excavación y cuna de arena). El tendido seguirá los lineamientos de la NCh399 y NCh2282/2. Tuberías y accesorios se unirán por soldadura química e irán enterrados sobre cuna de arena a una profundidad mínima según se detalla en los planos. Tuberías en \textbf{Metro lineal (ml)}, válvulas por \textbf{Unidad (un)}, fitting en \textbf{Global (gl)} y tierra en \textbf{Metro cúbico (m³)}.

\subsection{Recolección e impulsión de aguas agroquímicas}

\subsubsection*{Red PVC 25 mm, Fitting, Tablero y Movimiento de Tierras}

Comprende la línea de impulsión en PVC Hidráulico Ø 25 mm PN10 resistente a químicos, las piezas de acople, el tablero eléctrico de automatización para energización posterior, instalado a 1.5 m del suelo. La conducción se tenderá bajo norma NCh399 a una profundidad mínima de 0.60 m para protegerla del tránsito agrícola. El tablero cumplirá la normativa SEC para intemperie. Medición por \textbf{Metro lineal (ml)} para tubos, \textbf{Global (gl)} para fitting y tablero, y \textbf{Metro cúbico (m³)} en excavaciones.

\textbf{Nota:} los equipos de bombeo necesarios para el sector 1 serán provistos por el equipo técnico del proyecto. No considerarlos dentro del presupuesto.

\section{Sector 2: construcción de estación experimental (humedales)}

\subsection{Estructura para tanques de almacenamiento}

\subsubsection*{Fundaciones de Hormigón Armado (Zapatas, Enfierradura y Anclajes)}

Esta sección abarca la ejecución de las fundaciones para la torre elevada de tanques, incluyendo el vertido de hormigón grado G25, la instalación de mallas de enfierradura A630-420H (\(\varnothing 12\text{~mm}@15\text{~cm}\)) y la colocación de placas base de acero A36 (\(300 \times 300 \times 16\text{~mm}\)) con pernos de anclaje A325 de Ø 3/4". Todo el proceso respetará las normas NCh170, NCh430 y NCh204. El hormigonado se realizará sobre un emplantillado pobre de 5 cm, compactando mediante vibrado de inmersión y aplicando curado húmedo continuo por 7 días. Las enfierraduras dispondrán de calzos que garanticen un recubrimiento mínimo de 5 cm. Las placas base se posicionarán mediante plantillas rígidas y se nivelarán con mortero. Las zapatas se medirán en \textbf{Metro cúbico (m³)}, la enfierradura en \textbf{Metro lineal (ml)} y las placas base en \textbf{Metro cuadrado (m²)}.

\subsubsection*{Estructura Metálica Principal y Plataforma de Soporte}

Comprende la fabricación, soldadura y montaje de la estructura soporte compuesta por pilares cajón (\(150 \times 150 \times 6\text{~mm}\)), vigas principales (\(150 \times 150 \times 6\text{~mm}\)), vigas secundarias (\(150 \times 100 \times 3\text{~mm}\)) en acero A270ES, y una cubierta de tableros contrachapados fenólicos de pino radiata de 15 mm. Se regirá por la NCh427, AWS D1.1 y NCh1970. Las uniones soldadas a tope y filete serán continuas, ejecutadas por soldadores calificados y pulidas mecánicamente. Los tableros de madera se atornillarán a la estructura secundaria cada 20 cm e irán protegidos con un sellador antihumedad. Se aplicará inspección visual de soldaduras (VT). Los perfiles se medirán en \textbf{Metro lineal (ml)} y los tableros en \textbf{Metro cuadrado (m²)}.

\subsubsection*{Estructura Secundaria, Esquema de Pintura y Barandas}

Incluye los arriostramientos secundarios (\(15 \times 15 \times 6\text{~mm}\)), el tratamiento de protección superficial anticorrosivo/epóxico y las barandas perimetrales de seguridad en tubo Ø 30 mm. Las superficies metálicas se limpiarán antes de aplicar dos manos de imprimante anticorrosivo epóxico y esmalte de terminación para intemperie. Las barandas se instalarán a 1.0 m de altura con rodapié incorporado según la OGUC. Se medirá el espesor de película seca de pintura y se efectuarán pruebas de carga puntual sobre barandas. La medición será por \textbf{Metro lineal (ml)} en perfiles y barandas, y \textbf{Global (gl)} para la pintura.

\subsection{Sistema hidráulico de estructura de almacenamiento}

\subsubsection*{Estanques, Dosificadores, Válvulas y Redes de la Torre}

Abarca la instalación sobre la plataforma de madera de estanques verticales standard de 1.000 L (Ref. Bioplastic) para almacenamiento de agua y estanques dosificadores graduados de 200 L para entregas por pulso, junto con sus válvulas de bola PN10 de 50 mm, tuberías de PVC PN10 de 50 mm y accesorios inyectados. Los estanques serán opacos con protección UV e irán anclados mecánicamente para resistir vientos. La red de 25 mm y 50 mm correspondiente a las alimentaciones de llegada (recolección e impulsión de aguas del estero; recolección e impulsión de aguas agroquímicas; recolección e impulsión de aguas servidas) se fijará ordenadamente a la estructura mediante abrazaderas metálicas con goma amortiguadora. Se ejecutará prueba de llenado por 24 horas y prueba hidráulica de presión operacional. Los estanques y válvulas se medirán por \textbf{Unidad (un)}, la tubería por \textbf{Metro lineal (ml)} y los accesorios de forma \textbf{Global (gl)}.

\subsection{Construcción del humedal y sistema hidráulico}

\subsubsection*{\texorpdfstring{Obras Civiles Base (Solerillas, Losa \(\mathbf{e} = \mathbf{20}\text{~cm}\) y Sub-base)}{Obras Civiles Base (Solerillas, Losa \textbackslash mathbf\{e\} = \textbackslash mathbf\{20\}\textbackslash text\{\textasciitilde cm\} y Sub-base)}}

Consiste en la delimitación perimetral con solerillas de \(20 \times 100\text{~cm}\), sobre una losa de hormigón de 20 cm de espesor y una sub-base de grava estabilizada compactada. Se aplicará la norma NCh170. Tras perfilar y compactar el terreno, se colocará la sub-base de grava y se vaciará la losa de hormigón H20/G20, asentando las solerillas sobre mortero de pega. Se controlará la nivelación horizontal de la losa admitiendo una tolerancia máxima de \(\pm 5\text{~mm}\). Las solerillas se medirán por \textbf{Metro lineal (ml)}, mientras que la losa y la sub-base en \textbf{Metro cúbico (m³)}.

\subsubsection*{Estanques IBC para Celdas de Humedales}

Comprende el suministro, modificación y montaje de 18 contenedores industriales IBC de 1.000 Litros de PEAD con jaula galvanizada para conformar los humedales de flujo vertical (VFCW). Los contenedores se cortarán en su cara superior puliendo los bordes y se les perforará la base para acomodar pasamuros estancos de drenaje e inspección. Antes de introducir los medios filtrantes, cada estanque se someterá a llenado total para certificar cero filtraciones. La unidad de medida es por \textbf{Unidad (un)}.

\subsubsection*{Estratos Filtrantes (Grava 19 mm, Gravilla 12 mm y Arena Gruesa)}

Abarca el lavado exhaustivo y la colocación estratificada de los medios filtrantes granulares dentro de las celdas IBC: una capa inferior de drenaje en Grava de 19 mm (0,15 m), una capa intermedia de filtración en Arena Gruesa silícea (0,60 m) y una capa superior de distribución en Gravilla de 12 mm (0,15 m), según la lámina 2.2. Los materiales deberán cumplir con la NCh163 y poseer menos del 1\% de finos arcillosos. El vertido se hará manualmente por capas sin compactación mecánica, sólo apisonado, para preservar la porosidad teórica del diseño. Se realizarán ensayes de lavado e inspección de espesores por capa. Se medirá por \textbf{Metro cúbico (m³)}.

\subsubsection*{Plantación de Macrófitas (Typha spp.)}

Consiste en la selección, suministro e hincado de 72 rizomas/plántulas de la especie \emph{Typha spp.} (Totora) distribuidas a razón de 4 plantas por contenedor IBC. Se seleccionarán especímenes sanos con sistema radicular activo y brotes firmes. La plantación se hará de forma manual en la capa de arena e inmediatamente se iniciará un riego de saturación con agua limpia. Se evaluará la tasa de prendimiento a los 15 y 30 días. Se medirá por \textbf{Unidad (un)}.

\subsubsection*{Sistema de Alimentación, Drenaje y Distribución Hidráulica}

Comprende las tuberías principales de distribución en PVC 50 mm PN10, los ramales superiores perforados en PVC 32 mm, la red inferior de recolección en PVC Sanitario 50 mm, válvulas de corte de 32 mm y sus conectores. Los ramales superiores se perforarán uniformemente para asegurar una distribución por pulso homogénea sobre el sustrato. Se aplicarán pruebas de descarga en simultáneo para corroborar la uniformidad de distribución de caudal. Las tuberías se medirán en \textbf{Metro lineal (ml)}, las válvulas por \textbf{Unidad (un)} y los accesorios en \textbf{Global (gl)}.

\subsection{Recolección de efluentes tratados por sistema de humedal}

\subsubsection*{Desinfección, Acopio de Efluente y Obras de Interconexión}

Incluye el suministro de cámaras de cloración de 100 L con dosificador, fosas de acopio final de PEAD de 3.250 L, canalizaciones sanitarias en PVC de 50 mm, válvulas de 50 mm, excavaciones y cunas de grava. Las fosas responderán a la NCh2038 y se instalarán bajo llenado hidrostático preventivo durante el relleno de tierras. La dosificación de cloro se ajustará para garantizar los estándares de desinfección en el efluente final. Las cámaras y fosas se medirán por \textbf{Unidad (un)}, tuberías en \textbf{Metro lineal (ml)} y tierras en \textbf{Metro cúbico (m³)}.

\section{Sector 3: sistema de riego y otros}

\subsection{Sistema de conducción de riego}

\subsubsection*{Red Matriz, Goteros, Válvulas y Preparación del Terreno}

Comprende la instalación de la red matriz de PVC Hidráulico 50 mm PN10 y las líneas laterales de manguera de PEAD Plansa Ø 16 mm equipadas con goteros autocompensados de \(2\text{~l/h}\) a \(50\text{~cm}\), distribuidas a lo largo de las filas de los 30 árboles de cerezo (organizados en 3 sectores experimentales de 10 plantas cada uno. Las tuberías y mangueras se medirán en \textbf{Metro lineal (ml)}, válvulas por \textbf{Unidad (un)}, fitting en \textbf{Global (gl)} y movimientos de tierra en \textbf{Metro cúbico (m³)}.

\subsection{Otros}

\subsubsection*{\texorpdfstring{Cierre Perimetral de Madera y Malla (\(\mathbf{H} = \mathbf{1},\mathbf{5}\text{~m}\))}{Cierre Perimetral de Madera y Malla (\textbackslash mathbf\{H\} = \textbackslash mathbf\{1\},\textbackslash mathbf\{5\}\textbackslash text\{\textasciitilde m\})}}

Abarca la construcción de un cerramiento de \(1,5\text{~m}\) de altura constituido por polines de madera impregnada CCA de Ø 4-5" distanciados a \(2,5\text{~m}\) y empotrados \(0,50\text{~m}\) en dados de hormigón, sobre los cuales se tensionará una malla electrosoldada galvanizada de \(1,5\text{~m}\) fijada mediante grapas de acero. Se verificará la alineación, aplomado y rigidez mecánica del conjunto. La unidad de medida es por \textbf{Metro lineal (ml)}
